The argument is powerful but it is not free. It rests on four assumptions. First, (a)~buyers know
that sellers possess private information about quality; if a seller may simply not know, silence
no longer implies unwillingness to reveal. Second, (b)~disclosure is costless, or at least cheap
relative to the price consequences of silence; a large literature examines how disclosure costs
truncate the unraveling process and leave a residual pool of non-disclosers
\citep{verrecchia_1983, jovanovic_1982}. Third, (c)~disclosure is credible, so that a claim about
quality cannot be fabricated; without verifiability, disclosure degenerates into cheap talk and
the unraveling logic never starts \citep{crawford_sobel_1982, stocken_credibility_2000}. Fourth, (d)~buyers are homogeneous in how much they value the
disclosed attribute, so that the price penalty for silence is felt by every seller with equal
force.
